\begin{frame}{Como identificar a frequência de corte}
  \small
  \begin{block}{Resposta em regime permanente senoidal}
    Uma entrada senoidal de amplitude $A$ produz uma saída de amplitude
    $A|H(j\omega)|$, com defasagem $\arg H(j\omega)$.
    Para $H(s)=K/(1+s\tau)$, o corte ocorre em $\omega_c=1/\tau$:
    \[
      |H(j\omega_c)|=\frac{K}{\sqrt{2}},\qquad \phi(\omega_c)=-45^\circ.
    \]
  \end{block}

  {\centering
    \begin{tabular}{lccc}
      \toprule
      Projeto nominal & Ganho DC & Ganho no corte & Corte em dB absolutos \\
      \midrule
      Sem carga / buffer & $1$ & $0{,}707$ & $-3{,}01\,\mathrm{dB}$ \\
      Carga direta & $0{,}4$ & $0{,}283$ & $-10{,}97\,\mathrm{dB}$ \\
      \bottomrule
    \end{tabular}\par
  }

  \begin{alertblock}{Referência da medição}
    A queda é de $3\,\mathrm{dB}$ em relação ao patamar de cada circuito.
    Usam-se amplitudes de entrada e saída com a mesma convenção, por exemplo $V_{pp}$.
  \end{alertblock}
\end{frame}
